\chapter{}
\label{chap:1}

\section{Умова}
Знайти афінні статистичні аналоги наведених нижче функцій. Обчислити нелінійність кожної з цих функцій.
\begin{enumerate}[label=\alph*)]
    \item (1, 0, 1, 1, 0, 0, 1, 0);
    \item (0, 0, 1, 1, 0, 0, 1, 1).
\end{enumerate}

\section{Розв'язання}

Спершу я виписав для себе теоретичні нюанси з лекцій.

\noindent Для булевої функції $f \in B_n$ коефіцієнти Уолша--Адамара обраховуються формулою:
\begin{equation*}
    \hat{f}(\alpha) = \sum\limits_{x \in V_n} (-1)^{f(x) \oplus \alpha x}, \qquad \alpha \in V_n.
\end{equation*}
Але ми використовуватимемо їх зв'язок із коефіцієнтами Фур'є $C_f(\alpha)$:
\begin{equation}
    \label{eq:AdamarCoef}
    \hat{f}(\alpha) = \begin{cases}
        2^n - 2\,C_f(0), & \alpha = 0,    \\
        -2\,C_f(\alpha), & \alpha \neq 0.
    \end{cases}
\end{equation}
Афінні статистичні аналоги визначаються ненульовими значеннями $\hat{f}(\alpha)$. Згідно з Твердженням~3.7
(ст. 133), якщо $c \in \{0,1\}$ та $X$ -- рівноймовірний вектор із $V_n$, то якість аналога:
\begin{equation}
    \label{eq:probability}
    \Pr\big(f(X) = \alpha_1 X_1 \oplus \ldots \oplus \alpha_n X_n \oplus c\big)
    = 1/2 \cdot \big(1 + (-1)^c \cdot 2^{-n}\hat{f}(\alpha)\big) = 1/2 \cdot (1 + 2^{-n}|\hat{f}(\alpha)|).
\end{equation}
Тому при $\hat{f}(\alpha) > 0$ аналогом є $l_{\alpha,0}(x) = \alpha x$, а при $\hat{f}(\alpha) < 0$, аналог
$l_{\alpha,1}(x) = \alpha x \oplus 1$ (Означення~3.2).

Означення нелінійності функції пов'язане з відстанню Геммінга вiд неї до класу афінних функцій: (Твердження~3.8 з
лекції):
\begin{equation*}
    N_f = 2^{n-1}\big(1 - 2^{-n}\max\limits_{\alpha \in V_n} |\hat{f}(\alpha)|\big)
\end{equation*}

\subsection*{а) $f = (1, 0, 1, 1, 0, 0, 1, 0)$}

Першим кроком, як не крути, треба обчислити коефіцієнти Фур'є. Найкраще це робити за допомогою швидкого
перетворення Адамара.

\newpage
\begin{table}[htpb]
    \centering
    \begin{tblr}{
        colspec = {c | c | c | c},
        row{1} = {font=\bfseries},
        hline{2,10} = {solid}, % pivot lines
        hline{6} = {1}{solid}, % first column
        hline{4,6,8} = {2}{solid}, % second column
        hline{2,3,4,5,6,7,8,9} = {3}{solid} % third column
            }
        $f$ & Крок 1 & Крок 2 & $C_f$ \\
        1   & 1      & 3      & 4     \\
        0   & 0      & 1      & 2     \\
        1   & 2      & $-1$   & $-2$  \\
        1   & 1      & $-1$   & 0     \\
        0   & 1      & 1      & 2     \\
        0   & 0      & 1      & 0     \\
        1   & 0      & 1      & 0     \\
        0   & 1      & $-1$   & 2     \\
    \end{tblr}
\end{table}
Традиційно, рівністю Парсеваля перевіримо правильність обрахунку:
\begin{equation*}
    \frac{1}{2^3}\sum C_f(\alpha)^2 = \frac{16+4+4+0+4+0+0+4}{8} = \frac{32}{8} = 4 = \operatorname{wt}(f)
\end{equation*}
Далі обраховуємо коефіцієнти Уолша-Адамара за формулою~\eqref{eq:AdamarCoef}.

$\hat{f}(0,0,0) = 2^3 - 2 \cdot C_f(0,0,0) = 8 - 8 = 0$. Для всіх інших ($\alpha \neq 0$) маємо:
$\hat{f}(\alpha) = -2 \cdot C_f(\alpha)$.
\begin{center}
    \renewcommand{\arraystretch}{1.2}
    \begin{tabular}{c|cccccccc}
        $\alpha$          & $000$ & $001$ & $010$ & $011$ & $100$ & $101$ & $110$ & $111$ \\
        \hline
        $C_f(\alpha)$     & $4$   & $2$   & $-2$  & $0$   & $2$   & $0$   & $0$   & $2$   \\
        $\hat{f}(\alpha)$ & $0$   & $-4$  & $4$   & $0$   & $-4$  & $0$   & $0$   & $-4$
    \end{tabular}
\end{center}
Так само, перевірка Парсевалем:
\begin{equation*}
    \sum \hat{f}(\alpha)^2 = 0 + 16 + 16 + 0 + 16 + 0 + 0 + 16 = 64 = 2^{2 \cdot n}, \text{ де } n=3.
\end{equation*}
Далі лишається порахувати афінні статистичні аналоги. Їх визначають ненульові значення $\hat{f}(\alpha)$. В нас їх 4.
У випадках, якщо $\hat{f}(\alpha) < 0$ -- доксорюємо ще одиницю $\oplus\, 1$:

\begin{center}
    \renewcommand{\arraystretch}{1.3}
    \begin{tabular}{c|c|c|c}
        $\alpha$  & $\hat{f}(\alpha)$ & аналог $l_{\alpha,c}$                & $\Pr(f = l_{\alpha,c})$ формула~\eqref{eq:probability} \\
        \hline
        $(0,0,1)$ & $-4$              & $x_3 \oplus 1$                       & 3/4                                                    \\
        $(0,1,0)$ & $4$               & $x_2$                                & 3/4                                                    \\
        $(1,0,0)$ & $-4$              & $x_1 \oplus 1$                       & 3/4                                                    \\
        $(1,1,1)$ & $-4$              & $x_1 \oplus x_2 \oplus x_3 \oplus 1$ & 3/4                                                    \\
    \end{tabular}
\end{center}
Усі чотири аналоги мають однакову якість рівну $0.75$. Ну і на останок, обчислюємо нелінійність функції $f$,
a.k.a відстань Геммінга до класу афінних функцій складає:
\begin{equation*}
    N_f = 2^{3-1}\big(1 - 2^{-3}\max_\alpha |\hat{f}(\alpha)|\big) = 4\big(1 - \tfrac{4}{8}\big) = 4 \cdot \tfrac{1}{2} = 2.
\end{equation*}

\newpage % FORCED
\subsection*{б) $f = (0, 0, 1, 1, 0, 0, 1, 1)$}

Фан факт: в даному пункті, значення $f$ залежить лише від $x_2$ (судячи з таблиці значень~будемо мати $f = x_2$).

Так само спершу обраховуємо коефіцієнти Фур'є, тут все дуже просто:
\begin{table}[htpb]
    \centering
    \begin{tblr}{
        colspec = {c | c | c | c},
        row{1} = {font=\bfseries},
        hline{2,10} = {solid}, % pivot lines
        hline{6} = {1}{solid}, % first column
        hline{4,6,8} = {2}{solid}, % second column
        hline{2,3,4,5,6,7,8,9} = {3}{solid} % third column
            }
        $f$ & Крок 1 & Крок 2 & $C_f$ \\
        0   & 0      & 2      & 4     \\
        0   & 0      & 2      & 0     \\
        1   & 2      & $-2$   & $-4$  \\
        1   & 2      & $-2$   & 0     \\
        0   & 0      & 0      & 0     \\
        0   & 0      & 0      & 0     \\
        1   & 0      & 0      & 0     \\
        1   & 0      & 0      & 0     \\
    \end{tblr}
\end{table}

\noindent Парсеваль підтверджує правильність:
\begin{equation*}
    \frac{1}{8}(16 + 0 + 16 + 0 + 0 + 0 + 0 + 0) = 4 = \operatorname{wt}(f).
\end{equation*}
Коефіцієнти Уолша-Адамара, обраховані за формулою~\eqref{eq:AdamarCoef}, є наступними:
\begin{center}
    \renewcommand{\arraystretch}{1.2}
    \begin{tabular}{c|cccccccc}
        $\alpha$          & $000$ & $001$ & $010$ & $011$ & $100$ & $101$ & $110$ & $111$ \\
        \hline
        $C_f(\alpha)$     & $4$   & $0$   & $-4$  & $0$   & $0$   & $0$   & $0$   & $0$   \\
        $\hat{f}(\alpha)$ & $0$   & $0$   & $8$   & $0$   & $0$   & $0$   & $0$   & $0$
    \end{tabular}
\end{center}
Теж Персивалем перевіряємо:
\begin{equation*}
    \sum \hat{f}(\alpha)^2 = 64 = 2^{2n}, \text{ де } n=3.
\end{equation*}
І третім кроком рахуємо афінні статистичні аналоги і ймовірності. Оскільки єдиний ненульовий коефіцієнт: 
$\hat{f}(0,1,0) = 8 > 0$, то і один статистичний аналог $l_{(0,1,0),\,0} = x_2$. Якість такого аналога:
\begin{equation*}
    \Pr(f = l_{\alpha,c}) = \frac{1}{2}(1 + 2^{-3}\cdot 8) = \frac{1}{2}(1 + 1) = 1.
\end{equation*}
Отримали одиницю, а це означає, що $f = x_2$ є тотожньою, тобто $f$ сама по собі є афінною функцією, а тому 
нелінійність функції $f$ складатиме очевидно що нуль:
\begin{equation*}
    N_f = 2^{3-1}\big(1 - 2^{-3}\cdot 8\big) = 4\big(1 - 1\big) = 0.
\end{equation*}
